% ============================================================
%  RAPID — Comprehensive Technical Analysis Report
%  LaTeX Source
%  Compile with: pdflatex → bibtex → pdflatex → pdflatex
%  or simply:    pdflatex RAPID_ANALYSIS.tex  (twice for TOC)
% ============================================================
\documentclass[12pt, a4paper, oneside]{report}

% ─── Packages ────────────────────────────────────────────────
\usepackage[T1]{fontenc}
\usepackage[utf8]{inputenc}
\usepackage{lmodern}
\usepackage[margin=1in, top=1.1in, bottom=1.1in]{geometry}
\usepackage{microtype}

% Typography
\usepackage{palatino}          % Main serif font
\usepackage{helvet}            % Sans-serif for headings
\usepackage{courier}           % Monospace

% Colors
\usepackage[dvipsnames,table]{xcolor}
\definecolor{rapidblue}{RGB}{30, 90, 160}
\definecolor{rapidgray}{RGB}{80, 80, 90}
\definecolor{rapidlight}{RGB}{235, 242, 252}
\definecolor{rapidaccent}{RGB}{220, 50, 50}
\definecolor{codebg}{RGB}{245, 245, 248}
\definecolor{progreen}{RGB}{0, 110, 55}
\definecolor{conred}{RGB}{160, 30, 30}

% Links & PDF metadata
\usepackage[
  colorlinks=true,
  linkcolor=rapidblue,
  urlcolor=rapidblue,
  citecolor=rapidblue,
  pdftitle={RAPID — Comprehensive Technical Analysis},
  pdfauthor={Project Analysis Agent},
  pdfsubject={RAG Application for Private Instant Deployment},
  bookmarks=true,
  bookmarksopen=true
]{hyperref}

% Headers / Footers
\usepackage{fancyhdr}
\usepackage{emptypage}

% Tables
\usepackage{booktabs}
\usepackage{tabularx}
\usepackage{longtable}
\usepackage{multirow}
\usepackage{array}
\usepackage{makecell}
\renewcommand\theadfont{\bfseries}

% Lists
\usepackage{enumitem}
\setlist[itemize]{leftmargin=1.5em, itemsep=2pt, topsep=4pt}
\setlist[enumerate]{leftmargin=1.5em, itemsep=2pt, topsep=4pt}
\setlist[description]{leftmargin=1.5em, itemsep=2pt, topsep=4pt}

% Code listings
\usepackage{listings}
\lstset{
  basicstyle=\ttfamily\footnotesize,
  backgroundcolor=\color{codebg},
  frame=single,
  framerule=0.4pt,
  rulecolor=\color{rapidgray!40},
  breaklines=true,
  breakatwhitespace=true,
  tabsize=2,
  showstringspaces=false,
  keywordstyle=\color{rapidblue}\bfseries,
  commentstyle=\color{rapidgray}\itshape,
  stringstyle=\color{progreen},
  numbers=none,
  xleftmargin=6pt,
  xrightmargin=6pt,
  aboveskip=8pt,
  belowskip=8pt,
}

% Callout boxes
\usepackage{tcolorbox}
\tcbuselibrary{skins, breakable, listings}

\newtcolorbox{infobox}[1][]{
  colback=rapidlight, colframe=rapidblue,
  fonttitle=\bfseries\sffamily, title=#1,
  arc=3pt, boxrule=0.8pt, left=6pt, right=6pt,
  top=4pt, bottom=4pt, breakable
}

\newtcolorbox{probox}[1][]{
  colback=green!5, colframe=progreen,
  fonttitle=\bfseries\sffamily\color{progreen}, title=#1,
  arc=3pt, boxrule=0.8pt, left=6pt, right=6pt,
  top=4pt, bottom=4pt, breakable
}

\newtcolorbox{conbox}[1][]{
  colback=red!4, colframe=conred,
  fonttitle=\bfseries\sffamily\color{conred}, title=#1,
  arc=3pt, boxrule=0.8pt, left=6pt, right=6pt,
  top=4pt, bottom=4pt, breakable
}

\newtcolorbox{codebox}[1][]{
  colback=codebg, colframe=rapidgray!30,
  fonttitle=\bfseries\sffamily\footnotesize, title=#1,
  arc=2pt, boxrule=0.5pt, left=4pt, right=4pt,
  top=2pt, bottom=2pt, breakable,
  fontupper=\ttfamily\footnotesize
}

% Section formatting
\usepackage{titlesec}
\titleformat{\chapter}[display]
  {\normalfont\huge\bfseries\sffamily\color{rapidblue}}
  {\chaptertitlename\ \thechapter}{12pt}{\Huge}
\titleformat{\section}
  {\normalfont\Large\bfseries\sffamily\color{rapidblue}}
  {\thesection}{0.8em}{}[\vspace{-4pt}\rule{\linewidth}{0.5pt}]
\titleformat{\subsection}
  {\normalfont\large\bfseries\sffamily\color{rapidgray}}
  {\thesubsection}{0.7em}{}
\titleformat{\subsubsection}
  {\normalfont\normalsize\bfseries\sffamily}
  {\thesubsubsection}{0.6em}{}

\titlespacing*{\chapter}{0pt}{0pt}{20pt}
\titlespacing*{\section}{0pt}{16pt}{6pt}
\titlespacing*{\subsection}{0pt}{12pt}{4pt}

% TOC depth
\setcounter{tocdepth}{2}
\setcounter{secnumdepth}{3}

% Misc
\usepackage{graphicx}
\usepackage{float}
\usepackage{caption}
\captionsetup{font=small, labelfont=bf, labelsep=period}
\usepackage{parskip}
\setlength{\parindent}{0pt}
\setlength{\parskip}{6pt}
\usepackage{soul}        % \hl{} highlight
\usepackage{amsmath}

% Row color for tables
\rowcolors{2}{rapidlight!40}{white}

% ─── Page style ───────────────────────────────────────────────
\pagestyle{fancy}
\fancyhf{}
\fancyhead[L]{\small\sffamily\color{rapidgray} RAPID — Technical Analysis}
\fancyhead[R]{\small\sffamily\color{rapidgray} \leftmark}
\fancyfoot[C]{\small\sffamily\color{rapidgray} \thepage}
\renewcommand{\headrulewidth}{0.4pt}
\renewcommand{\headrule}{\color{rapidblue!50}\hrule width\headwidth height\headrulewidth}

% ─── Custom commands ──────────────────────────────────────────
\newcommand{\feature}[1]{\subsubsection*{\color{rapidblue}\sffamily #1}}
\newcommand{\kw}[1]{\texttt{\small #1}}
\newcommand{\pro}[1]{\textcolor{progreen}{\textbf{+}} #1}
\newcommand{\con}[1]{\textcolor{conred}{\textbf{--}} #1}
\newcommand{\rapidsys}{\textbf{\color{rapidblue}RAPID}}

% ─── Document ─────────────────────────────────────────────────
\begin{document}

% ═══════════════════════════════════════════════════════════════
%  COVER PAGE
% ═══════════════════════════════════════════════════════════════
\begin{titlepage}
  \pagecolor{rapidblue}
  \color{white}
  \vspace*{2cm}
  \begin{center}
    {\fontsize{48}{54}\selectfont\bfseries\sffamily RAPID}\\[8pt]
    {\Large\sffamily RAG Application for Private Instant Deployment}\\[2cm]
    \rule{0.6\linewidth}{0.8pt}\\[1.5cm]
    {\LARGE\sffamily\bfseries Comprehensive Technical Analysis}\\[0.5cm]
    {\large\sffamily System Design, Feature Evaluation, Workflows,}\\
    {\large\sffamily Architecture \& Technology Stack}\\[2.5cm]
    \rule{0.6\linewidth}{0.4pt}\\[1cm]
    {\large\sffamily Analysis Date: February 27, 2026}\\[0.4cm]
    {\large\sffamily Document Version: 1.0}\\[2cm]
    {\normalsize\sffamily\color{white!80}
      Designed \& Built by \textbf{Ayush Chhoker}\\[0.3cm]
      Individual Project
    }
  \end{center}
  \vfill
\end{titlepage}
\nopagecolor

% ═══════════════════════════════════════════════════════════════
%  TABLE OF CONTENTS
% ═══════════════════════════════════════════════════════════════
\tableofcontents
\newpage

% ═══════════════════════════════════════════════════════════════
%  CHAPTER 1 — EXECUTIVE SUMMARY
% ═══════════════════════════════════════════════════════════════
\chapter{Executive Summary}

\rapidsys{} is an enterprise-grade \textbf{private AI assistant platform} that enables any
organization to query its own internal documents and databases using plain English — without
transmitting sensitive data to third-party cloud services.

Think of it as a private version of ChatGPT that ``knows'' everything your organization has
ever written down. Upload a legal contract, a financial report, or a CSV of last year's
sales data, and \rapidsys{} automatically determines the best way to answer questions about
it — either through intelligent document search (RAG) or through auto-generated SQL queries
on structured data.

\section*{What Makes RAPID Different}

Most AI document tools require manual configuration of chunk sizes, embedding models, and
search strategies. \rapidsys{} eliminates this complexity with an \textbf{Intelligent
Auto-RAG} system: it automatically detects what kind of document was uploaded (legal
contract, academic paper, spreadsheet, etc.) and configures the entire pipeline with
research-backed optimal parameters, invisible to the end user.

\section*{Key Capabilities at a Glance}

\begin{table}[H]
\centering
\rowcolors{2}{rapidlight!50}{white}
\begin{tabularx}{\linewidth}{>{\bfseries}lX}
\toprule
\rowcolor{rapidblue!10}
Capability & Details \\
\midrule
Document types & 32+ file types (PDF, Word, Excel, CSV, code, JSON, YAML, and more) \\
LLM providers & OpenAI, Anthropic Claude, OpenRouter, Ollama (local), LM Studio (local) \\
Embedding providers & SentenceTransformers (local), OpenAI, Ollama, HuggingFace \\
Search modes & Semantic (vector), Keyword (BM25), Hybrid (RRF fusion) \\
Data query & Natural Language \(\rightarrow\) SQL for spreadsheets and databases \\
Multi-tenancy & Organizations, departments, groups with fine-grained permissions \\
Privacy & Fully self-hostable; no data leaves your infrastructure \\
\bottomrule
\end{tabularx}
\caption{RAPID — Key Capabilities}
\end{table}

% ═══════════════════════════════════════════════════════════════
%  CHAPTER 2 — PROJECT OVERVIEW
% ═══════════════════════════════════════════════════════════════
\chapter{Project Overview \& Purpose}

\section{What Is RAPID?}

\rapidsys{} stands for \textbf{R}AG \textbf{A}pplication for \textbf{P}rivate \textbf{I}nstant
\textbf{D}eployment. It is a full-stack AI application that serves as a private, organizational
knowledge-base assistant. Users can:

\begin{itemize}
  \item Upload documents in 32+ formats
  \item Ask questions in plain English
  \item Get answers sourced directly from their uploaded documents
  \item Query spreadsheet data without knowing SQL
  \item Connect to live databases (PostgreSQL, MySQL, Snowflake, BigQuery)
  \item Get web-augmented answers when internal documents lack the information
\end{itemize}

\section{The Problem It Solves}

Organizations accumulate vast amounts of knowledge locked in files: contracts, reports,
policies, research papers, spreadsheets. Employees spend hours searching through these files
manually. Traditional search tools find documents but do not \emph{answer questions}.

Existing commercial AI tools (ChatGPT Enterprise, Google Gemini for Workspace) have two
critical drawbacks for many organizations:
\begin{enumerate}
  \item \textbf{Data privacy}: documents are sent to third-party servers.
  \item \textbf{Configuration burden}: optimal AI configuration requires expertise most
        organizations do not have.
\end{enumerate}

\rapidsys{} solves both: it runs entirely on your infrastructure (no data leaves), and it
automatically configures itself based on what you upload.

\section{Target Audience}

\begin{description}
  \item[Primary:] Organizations with sensitive data that cannot leave their environment —
        legal firms, healthcare providers, financial institutions, government agencies,
        research institutions.
  \item[Secondary:] Any organization wanting a self-hosted AI knowledge base without needing
        a team of ML engineers to configure it.
\end{description}

\section{Core Value Proposition}

\begin{infobox}[Core Value]
``Upload your files, ask your questions, get answers — automatically and privately.''
\end{infobox}

RAPID's value comes from three pillars:
\begin{enumerate}
  \item \textbf{Zero-configuration intelligence} — Auto-detects document type and auto-configures
        the entire AI pipeline.
  \item \textbf{Privacy by design} — Deployable entirely on-premises with local LLMs (Ollama)
        and local embeddings (SentenceTransformers).
  \item \textbf{Dual-pipeline routing} — Automatically routes structured data to SQL and
        unstructured text to RAG.
\end{enumerate}

% ═══════════════════════════════════════════════════════════════
%  CHAPTER 3 — HOW IT WORKS
% ═══════════════════════════════════════════════════════════════
\chapter{How It Works — Deep Dive}

\section{Architecture Overview}

\rapidsys{} uses a \textbf{three-tier layered architecture}:

\begin{codebox}[System Architecture]
PRESENTATION TIER
  Streamlit Web UI  <-->  FastAPI REST Backend
              |
INTELLIGENCE TIER
  DocumentClassifier --> AutoConfigService
  ChunkingOptimizer  --> TextPreprocessor
  QueryDecomposer    --> MultiAgentOrchestrator
  ConfidenceScorer   --> SpeculativePipeline
              |
STORAGE TIER
  ChromaDB (vectors)    SQLite (users, metadata)
  BM25 Index (keywords) SQLite (tabular data)
  NetworkX (knowledge graph)
\end{codebox}

\section{The Upload Flow — How Documents Are Processed}

When a user uploads a file, the following steps happen automatically:

\begin{enumerate}
  \item \textbf{File Validation} — \kw{FileValidator} checks extension and MIME type against
        a whitelist of 32+ supported formats. Files exceeding the size limit are rejected.

  \item \textbf{Document Classification} — \kw{DocumentClassifier} analyzes the file using
        heuristics (no LLM call needed):
        \begin{itemize}
          \item Checks file extension (\kw{.csv}, \kw{.parquet} \(\rightarrow\) always tabular)
          \item For text files, reads a sample and scores keyword frequency against domain
                vocabularies (legal terms, medical terms, academic terms, etc.)
          \item Outputs: \kw{doc\_type} (e.g., ``academic'', ``legal'', ``tabular'') and
                \kw{doc\_subtype} (e.g., ``research\_paper'', ``contract'',
                ``financial\_spreadsheet'')
        \end{itemize}

  \item \textbf{Pipeline Routing} — Based on classification:
        \begin{itemize}
          \item Tabular files \(\rightarrow\) \kw{TabularPipeline} (loads into SQLite, enables NL-to-SQL)
          \item All other files \(\rightarrow\) \kw{RAGEngine} (chunks, embeds, stores in ChromaDB)
        \end{itemize}

  \item \textbf{Tabular Pipeline} (for spreadsheets):
        \begin{enumerate}
          \item Loads CSV/Excel/Parquet into a SQLite table
          \item Registers as a database connection in \kw{CloudDatabaseService}
          \item Subsequent queries generate SQL automatically
        \end{enumerate}

  \item \textbf{RAG Pipeline} (for text documents):
        \begin{enumerate}
          \item \kw{AutoConfigService} looks up research-backed optimal settings
          \item \kw{TextPreprocessor} cleans OCR artifacts, normalizes whitespace
          \item \kw{LanguageDetector} checks language; switches to multilingual embedding if needed
          \item \kw{ChunkingOptimizer} selects from 11 strategies based on document type
          \item Text is split into chunks, embedded (converted to vectors), stored in ChromaDB
          \item BM25 keyword index is built for hybrid retrieval
          \item Knowledge graph entities extracted (optional, asynchronous)
          \item Document registered in \kw{DataCatalogService} with topic keywords
        \end{enumerate}
\end{enumerate}

\section{The Query Flow — How Questions Are Answered}

\begin{enumerate}
  \item \textbf{Query Decomposition} — \kw{QueryDecomposer} checks if the query is complex
        (multiple questions, comparisons, conjunctions). If so, it splits it into focused
        sub-queries.

  \item \textbf{Intent Classification} — The \kw{MultiAgentOrchestrator} (powered by
        LangGraph) classifies query intent: \kw{general\_chat}, \kw{query\_agent},
        \kw{db\_agent}, \kw{multi\_source}, or \kw{graph\_query}.

  \item \textbf{Retrieval} — For RAG queries:
        \begin{itemize}
          \item Semantic search in ChromaDB (vector similarity)
          \item Keyword search via BM25 index
          \item Results merged using Reciprocal Rank Fusion (RRF)
          \item RBAC filter applied — user only sees documents they can access
        \end{itemize}

  \item \textbf{Generation} — LLM generates an answer from retrieved context. Optionally uses
        the Speculative Pipeline (draft fast \(\rightarrow\) verify with larger model).

  \item \textbf{Confidence Scoring} — \kw{ConfidenceScorer} evaluates on three dimensions:
        \begin{itemize}
          \item Context Relevance (30 \%) — Are chunks actually about the question?
          \item Faithfulness (50 \%) — Is the answer grounded in retrieved text?
          \item Completeness (20 \%) — Did the answer address what was asked?
        \end{itemize}
        Score $< 0.40$: retry with different strategy. Score $0.40$--$0.65$: low-confidence
        flag. Score $\geq 0.65$: high confidence.

  \item \textbf{Response \& Citation} — Answer returned with source citations. Conflicts
        between sources are detected and flagged.
\end{enumerate}

\section{Data Flow — Text Diagram}

\begin{codebox}[End-to-End Query Pipeline]
User Question
    |
    v
QueryDecomposer ----complex?----> split into sub-queries
    |
    v
MultiAgentOrchestrator (LangGraph state machine)
    |
    |-- classify_node: intent label
    |
    |-- [general_chat]  --> direct_answer_node --> LLM --> response
    |
    |-- [query_agent]   --> two_stage_query:
    |       |-- semantic search (ChromaDB)
    |       |-- keyword search (BM25)
    |       +-- RRF fusion --> top chunks --> LLM generates answer
    |
    |-- [db_agent]      --> schema context --> LLM --> SQL --> execute --> format
    |
    +-- [multi_source]  --> both paths above --> fusion_node --> merged answer
    |
    v
verify_node (ConfidenceScorer)
    |
    |-- score >= 0.65  --> return answer
    |-- score 0.40-0.65 --> return with low-confidence flag
    +-- score < 0.40   --> repair_node (HyDE) --> retry once
                                    |
                               still low --> partial_deliver_node
    |
    v
Response with sources, confidence level, conflict warnings
\end{codebox}

% ═══════════════════════════════════════════════════════════════
%  CHAPTER 4 — FEATURE ANALYSIS
% ═══════════════════════════════════════════════════════════════
\chapter{Comprehensive Feature Analysis}

\section{Feature 1: Intelligent Auto-RAG}

\subsection*{Description}
The flagship feature. \kw{DocumentClassifier} identifies document type; \kw{AutoConfigService}
maps that to research-backed optimal chunking and retrieval settings — all invisible to users.

\subsection*{Research-Backed Configurations}

\begin{table}[H]
\centering
\small
\rowcolors{2}{rapidlight!40}{white}
\begin{tabularx}{\linewidth}{lXllX}
\toprule
\rowcolor{rapidblue!10}
\textbf{Document Type} & \textbf{Chunk} & \textbf{Overlap} & \textbf{Mode} & \textbf{Benchmark} \\
\midrule
FAQ & 220 words & 25 & Keyword & NQ / TriviaQA \\
Policy / Procedure & 380 words & 50 & Hybrid & SQuAD / FinQA \\
General narrative & 512 words & 64 & Hybrid & BEIR multi-domain \\
Academic paper & 950 words & 150 & Hybrid & QASPER (+12\% F1) \\
Legal contract & 780 words & 100 & Hybrid & CUAD benchmark \\
Medical / clinical & 500 words & 60 & Hybrid & PubMedQA / MedQA \\
Financial report & 600 words & 80 & Hybrid & FinanceBench \\
Source code & 300 words & 50 & Hybrid & Function-level \\
\bottomrule
\end{tabularx}
\caption{AutoConfigService — research-backed pipeline configurations}
\end{table}

\begin{probox}[Benefits]
\begin{itemize}
  \item Zero configuration for end users
  \item Demonstrably better retrieval quality (e.g., +12\% F1 on academic papers vs.\ 512-word baseline)
  \item Consistent, reproducible behavior based on published research
\end{itemize}
\end{probox}

\begin{conbox}[Drawbacks]
\begin{itemize}
  \item Heuristic classification can misclassify ambiguous documents
  \item Configuration table is static (no online learning or feedback loop)
  \item Admin cannot override individual document settings through the UI
  \item No A/B testing infrastructure to validate configurations against real data
\end{itemize}
\end{conbox}

% ─────────────────────────────────────────────────────────────
\section{Feature 2: Eleven-Strategy Chunking Optimizer}

\subsection*{Description}
Instead of a single fixed chunking approach, \rapidsys{} implements 11 different strategies
and automatically selects the best one for each document type.

\subsection*{Strategies}

\begin{enumerate}
  \item \kw{fixed\_word} — Fixed word-count windows (baseline)
  \item \kw{fixed\_sentence} — Group N sentences per chunk
  \item \kw{paragraph} — Split on blank lines
  \item \kw{semantic} — Embedding-based split at semantic discontinuities
  \item \kw{recursive} — LangChain-style recursive character splitting
  \item \kw{sliding\_window} — 50\% overlap for dense Q\&A content
  \item \kw{token\_aware} — Accounts for word-to-token ratio ($\times 1.3$)
  \item \kw{section\_aware} — Detects markdown headers, ALL-CAPS lines, numbered sections
  \item \kw{code\_aware} — Splits at function/class boundaries (Python, JS, Java, Go, Rust)
  \item \kw{table\_intact} — Keeps markdown/HTML tables as atomic units
  \item \kw{step\_aware} — Splits at numbered step boundaries for procedural documents
\end{enumerate}

\subsection*{Selection Logic}
\begin{itemize}
  \item Subtype overrides: \kw{faq} \(\to\) \kw{step\_aware}, \kw{markdown\_table} \(\to\) \kw{table\_intact}, \kw{research\_paper} \(\to\) \kw{section\_aware}
  \item Type defaults: \kw{code} \(\to\) \kw{code\_aware}, \kw{academic} \(\to\) \kw{section\_aware}, \kw{legal} \(\to\) \kw{paragraph}
  \item Optional trial scoring: runs top-2 candidates and picks the better one by heuristic quality score
\end{itemize}

\begin{probox}[Benefits]
\begin{itemize}
  \item Preserves document structure (tables stay together, code functions stay together)
  \item Better retrieval precision for structured documents
  \item No configuration required from users
\end{itemize}
\end{probox}

\begin{conbox}[Drawbacks]
\begin{itemize}
  \item Semantic strategy requires loading a SentenceTransformer model — adds latency
  \item Trial scoring doubles processing time for ambiguous documents
  \item Quality scoring is heuristic only, not based on actual retrieval outcomes
  \item \kw{code\_aware} detects Python/JS/Java patterns but may miss other languages
\end{itemize}
\end{conbox}

% ─────────────────────────────────────────────────────────────
\section{Feature 3: Dual-Pipeline — RAG + SQL Routing}

\subsection*{Description}
When a user uploads a CSV, Excel, or Parquet file, instead of embedding rows of numbers into
vectors (which performs poorly), \rapidsys{} loads the data into a local SQLite database and
routes all queries through an NL-to-SQL pipeline.

\subsection*{How It Works}
\begin{enumerate}
  \item \kw{TabularPipeline.ingest()} reads the file with automatic encoding/delimiter detection
  \item Data loaded into \kw{data/tabular\_uploads.db} as a named SQLite table
  \item Connection registered in \kw{CloudDatabaseService}
  \item On query, the orchestrator's \kw{DatabaseProxyAgent} gets schema context, LLM generates SQL
  \item SQL executed; results returned as natural language
\end{enumerate}

\begin{probox}[Benefits]
\begin{itemize}
  \item Handles datasets with hundreds of thousands of rows efficiently
  \item Exact numeric answers (vector search on numbers is unreliable)
  \item Supports CSV, Excel (multi-sheet), Parquet, and JSON arrays
  \item Multiple sheets become multiple tables with clear naming
\end{itemize}
\end{probox}

\begin{conbox}[Drawbacks]
\begin{itemize}
  \item SQL generation can fail for complex aggregations or ambiguous column names
  \item 500,000 row safety cap may truncate large datasets
  \item Excel files with merged cells or pivot tables may not load correctly
  \item LLM-generated SQL is not validated before execution (SQL injection risk via adversarial prompting)
\end{itemize}
\end{conbox}

% ─────────────────────────────────────────────────────────────
\section{Feature 4: Multi-Agent Orchestrator with LangGraph}

\subsection*{Description}
The query handling system is implemented as a directed state graph using LangGraph. This is a
CRAG-inspired (Corrective RAG) architecture.

\subsection*{Graph Nodes}

\begin{table}[H]
\centering
\small
\rowcolors{2}{rapidlight!40}{white}
\begin{tabularx}{\linewidth}{lXll}
\toprule
\rowcolor{rapidblue!10}
\textbf{Node} & \textbf{Responsibility} & \textbf{Input} & \textbf{LLM?} \\
\midrule
\kw{classify\_node} & Intent classification & query & Yes \\
\kw{direct\_answer\_node} & General chat, no retrieval & query + history & Yes \\
\kw{query\_agent\_node} & RAG retrieval + generation & query & Yes \\
\kw{db\_agent\_node} & SQL generation + execution & query + schema & Yes \\
\kw{multi\_source\_node} & Parallel RAG + SQL & query & Parallel \\
\kw{fusion\_node} & Merge multi-source results & rag + db & Yes \\
\kw{graph\_query\_node} & Knowledge graph traversal & query & Yes \\
\kw{verify\_node} & Confidence scoring & answer + chunks & Optional \\
\kw{repair\_node} & HyDE retry for low confidence & low answer & Yes \\
\kw{partial\_deliver\_node} & Return best available & any & No \\
\bottomrule
\end{tabularx}
\caption{LangGraph node responsibilities}
\end{table}

\begin{probox}[Benefits]
\begin{itemize}
  \item Clean separation of concerns between query types
  \item Automatic retry with different strategy for low-confidence answers
  \item Complex queries can use multiple sources simultaneously
  \item State is inspectable for debugging and tracing
\end{itemize}
\end{probox}

\begin{conbox}[Drawbacks]
\begin{itemize}
  \item LangGraph adds complexity; debugging state machine issues is non-trivial
  \item LLM classification step adds latency to every single query
  \item Repair mechanism (HyDE) makes another LLM call, doubling cost for low-confidence queries
  \item Concurrent parallel queries may hit provider rate limits
\end{itemize}
\end{conbox}

% ─────────────────────────────────────────────────────────────
\section{Feature 5: 3D Confidence Scoring System}

\subsection*{Description}
After every answer is generated, the \kw{ConfidenceScorer} evaluates it on three dimensions
and produces a structured verdict with retry recommendations.

\subsection*{Scoring Formula}
\[
  \text{overall} = 0.30 \times \text{ContextRelevance}
                 + 0.50 \times \text{Faithfulness}
                 + 0.20 \times \text{Completeness}
\]

\subsection*{Verdict Thresholds}

\begin{table}[H]
\centering
\rowcolors{2}{rapidlight!40}{white}
\begin{tabularx}{0.7\linewidth}{llX}
\toprule
\rowcolor{rapidblue!10}
\textbf{Range} & \textbf{Verdict} & \textbf{Action} \\
\midrule
$\geq 0.65$ & \textcolor{progreen}{\textbf{High}} & Return answer \\
$0.40$--$0.65$ & \textcolor{orange}{\textbf{Medium}} & Return with warning flag \\
$< 0.40$ & \textcolor{conred}{\textbf{Low}} & Trigger retry \\
\bottomrule
\end{tabularx}
\caption{Confidence verdict thresholds}
\end{table}

\subsection*{Additional Capabilities}
\begin{itemize}
  \item \textbf{Unanswerable Detection}: If $< 20\%$ of query keywords appear in chunks, flag as unanswerable.
  \item \textbf{Conflict Detection}: Cross-chunk numeric conflict detection (same subject, different values).
  \item \textbf{Citation-level Faithfulness}: Per-sentence support tracking with source chunk attribution.
\end{itemize}

\begin{probox}[Benefits]
\begin{itemize}
  \item Prevents hallucinated answers from being served confidently
  \item Conflict detection helps users know when sources disagree
  \item Sentence-level citation tracking for auditability
\end{itemize}
\end{probox}

\begin{conbox}[Drawbacks]
\begin{itemize}
  \item Heuristic faithfulness can give high scores to answers that merely repeat context verbatim
  \item LLM faithfulness check adds latency and cost per query
  \item Numeric conflict detection is pattern-based and generates false positives (different years' data)
  \item Completeness scoring penalizes valid short answers
\end{itemize}
\end{conbox}

% ─────────────────────────────────────────────────────────────
\section{Feature 6: Hybrid Search — Semantic + BM25 + RRF Fusion}

\subsection*{Description}
\rapidsys{} combines vector similarity search and keyword matching using Reciprocal Rank
Fusion (RRF) — a rank-aggregation technique robust to score scale differences.

\subsection*{Fusion Formula}
For a document $d$ appearing at rank $r_s$ in semantic results and rank $r_k$ in keyword results:
\[
  \text{RRF}(d) = \frac{\alpha}{r_s + k} + \frac{(1-\alpha)}{r_k + k}, \quad k = 60
\]

\subsection*{Search Mode Selection}
\begin{description}
  \item[\kw{keyword}:] FAQ / factual documents (BM25 only)
  \item[\kw{semantic}:] Structured / nested documents
  \item[\kw{hybrid}:] All other types (default)
\end{description}

\begin{probox}[Benefits]
\begin{itemize}
  \item Catches cases where semantic search fails (exact term match, specific numbers, abbreviations)
  \item Catches cases where keyword search fails (synonyms, paraphrasing)
  \item RRF is parameter-robust and does not require tuning score scales
\end{itemize}
\end{probox}

\begin{conbox}[Drawbacks]
\begin{itemize}
  \item BM25 index stored as JSON on disk; fully reloaded on restart (no incremental update)
  \item BM25 model rebuilt from scratch on every new document (O(n) in total corpus)
  \item No deduplication of overlapping chunks from both retrieval methods
  \item $\alpha = 0.5$ weighting is fixed; not tunable per document type
\end{itemize}
\end{conbox}

% ─────────────────────────────────────────────────────────────
\section{Feature 7: Knowledge Graph Builder}

\subsection*{Description}
During ingestion, \rapidsys{} can optionally extract entities and relationships using
LLM-based Named Entity Recognition (NER), building a NetworkX directed graph stored
per-user as a JSON file.

\begin{probox}[Benefits]
\begin{itemize}
  \item Enables relationship queries (``Who reports to the CFO?'' / ``Which departments use which systems?'')
  \item Per-user isolation preserves privacy
  \item NetworkX provides rich graph algorithms (shortest path, centrality, etc.)
\end{itemize}
\end{probox}

\begin{conbox}[Drawbacks]
\begin{itemize}
  \item LLM-based extraction is slow and costly for large document sets
  \item Entity resolution not implemented — ``Apple Inc.'' and ``Apple'' may become separate nodes
  \item JSON persistence does not scale for large graphs (loaded entirely into memory)
  \item No graph visualization in the UI
\end{itemize}
\end{conbox}

% ─────────────────────────────────────────────────────────────
\section{Feature 8: Multi-Provider LLM \& Embedding Support}

\subsection*{LLM Providers}

\begin{table}[H]
\centering
\rowcolors{2}{rapidlight!40}{white}
\begin{tabularx}{\linewidth}{llX}
\toprule
\rowcolor{rapidblue!10}
\textbf{Provider} & \textbf{Type} & \textbf{Notes} \\
\midrule
OpenAI & Cloud & GPT-3.5, GPT-4, GPT-4o \\
Anthropic & Cloud & Claude 3 Haiku / Sonnet / Opus \\
OpenRouter & Cloud & 100+ models via single API \\
Ollama & Local & Any pulled model (Llama, Mistral, etc.) \\
LM Studio & Local & OpenAI-compatible local server \\
\bottomrule
\end{tabularx}
\caption{Supported LLM providers}
\end{table}

\subsection*{Embedding Providers}

\begin{table}[H]
\centering
\rowcolors{2}{rapidlight!40}{white}
\begin{tabularx}{\linewidth}{llX}
\toprule
\rowcolor{rapidblue!10}
\textbf{Provider} & \textbf{Type} & \textbf{Notes} \\
\midrule
SentenceTransformers & Local & Default; \kw{all-MiniLM-L6-v2} (384d) \\
Ollama & Local & \kw{nomic-embed-text}, \kw{mxbai-embed-large} \\
OpenAI & Cloud & \kw{text-embedding-3-small/large} \\
HuggingFace & Cloud & Via Inference API \\
\bottomrule
\end{tabularx}
\caption{Supported embedding providers}
\end{table}

\begin{probox}[Benefits]
\begin{itemize}
  \item Organizations can start with local models (100\% private) and switch to cloud if needed
  \item Multilingual model auto-switching via LanguageDetector
  \item No vendor lock-in
\end{itemize}
\end{probox}

\begin{conbox}[Drawbacks]
\begin{itemize}
  \item LLM context windows differ across providers (not accounted for in chunk assembly)
  \item Model dimension mismatch if embedding provider is switched after documents are indexed
  \item Anthropic provides no embedding endpoints — only usable as LLM
\end{itemize}
\end{conbox}

% ─────────────────────────────────────────────────────────────
\section{Feature 9: Enterprise Security}

\subsection*{Authentication}
\begin{itemize}
  \item JWT tokens (HS256) with 30-minute expiration
  \item bcrypt password hashing with salt
  \item OAuth/OIDC SSO: Google, Microsoft, generic OIDC, Dropbox
  \item Rate limiting: in-memory sliding window (configurable limit/window)
  \item Password strength: min 8 chars, uppercase + lowercase + digit
\end{itemize}

\subsection*{Authorization (RBAC)}

\begin{table}[H]
\centering
\rowcolors{2}{rapidlight!40}{white}
\begin{tabularx}{\linewidth}{lX}
\toprule
\rowcolor{rapidblue!10}
\textbf{Layer} & \textbf{Details} \\
\midrule
Roles & \kw{admin}, \kw{manager}, \kw{user} \\
Document access levels & \kw{private}, \kw{group}, \kw{org}, \kw{public} \\
Permission check order & owner $\to$ admin $\to$ access level $\to$ allowed users $\to$ allowed roles $\to$ allowed groups \\
Audit log & Rotating file (10 MB max, 5 backups), JSON-structured \\
\bottomrule
\end{tabularx}
\caption{RBAC permission model}
\end{table}

\begin{conbox}[Drawbacks]
\begin{itemize}
  \item Rate limiting is \textbf{in-memory only} — resets on restart; does not work across multiple server instances
  \item JWT secret stored in a file, not a proper secrets manager
  \item No refresh token mechanism — users must re-login every 30 minutes
  \item SQLite does not scale well for high-concurrency writes
  \item Conversation messages stored in plaintext (not encrypted at rest)
\end{itemize}
\end{conbox}

% ─────────────────────────────────────────────────────────────
\section{Feature 10: Cloud Storage Connectors}

\rapidsys{} connects to five cloud storage services to pull documents directly from
organizational repositories.

\begin{table}[H]
\centering
\rowcolors{2}{rapidlight!40}{white}
\begin{tabularx}{\linewidth}{llX}
\toprule
\rowcolor{rapidblue!10}
\textbf{Service} & \textbf{Library} & \textbf{Auth Method} \\
\midrule
AWS S3 & \kw{boto3} & IAM credentials / env vars \\
Azure Blob Storage & \kw{azure-storage-blob} & Connection string / Azure AD \\
Google Drive & \kw{google-api-python-client} & OAuth 2.0 \\
Microsoft OneDrive & \kw{msal} + Graph API & OAuth 2.0 / MSAL \\
Dropbox & \kw{dropbox} SDK & OAuth 2.0 \\
\bottomrule
\end{tabularx}
\caption{Supported cloud storage connectors}
\end{table}

% ─────────────────────────────────────────────────────────────
\section{Feature 11: Web Search Augmentation}

When internal documents lack an answer (\kw{is\_unanswerable()} returns True), \rapidsys{}
optionally falls back to web search (disabled by default).

\textbf{Supported providers:} Tavily (AI-optimized), Google Custom Search, Bing Search API.

% ─────────────────────────────────────────────────────────────
\section{Feature 12: Speculative RAG Pipeline}

\begin{codebox}[Speculative RAG Flow]
Small LLM (fast) --> Draft answer
        |
        v
ConfidenceScorer: faithfulness >= threshold (0.70)?
        |
   YES  +-----> Return draft (fast path, lower cost)
        |
   NO   +-----> Large LLM corrects draft --> Return corrected answer
\end{codebox}

\begin{probox}[Benefits]
\begin{itemize}
  \item Lower cost: most queries answered by the draft (fast + cheap)
  \item Higher quality: verifier catches errors before serving
  \item Configurable threshold for quality/cost trade-off
\end{itemize}
\end{probox}

% ─────────────────────────────────────────────────────────────
\section{Feature 13: Query Decomposition}

Complex multi-part questions are automatically split into focused sub-queries.

\textbf{Example:}
\begin{quote}
``What was Q3 revenue \textit{and} who approved the budget?''\\
\(\Downarrow\)\\
Sub-query A: ``What was Q3 revenue?'' \(\to\) SQL pipeline\\
Sub-query B: ``Who approved the budget?'' \(\to\) RAG pipeline
\end{quote}

% ─────────────────────────────────────────────────────────────
\section{Feature 14: Data Catalog Service}

Every uploaded document is registered in a lightweight catalog with topic keywords extracted
by frequency analysis. At query time, the catalog is used to scope retrieval to only the
most relevant documents, avoiding irrelevant search.

\begin{codebox}[Catalog Schema]
data_catalog:
  doc_id, filename, username, doc_type, doc_subtype,
  pipeline, topics (JSON), stats (JSON),
  conn_id, doc_date, upload_time
\end{codebox}

% ─────────────────────────────────────────────────────────────
\section{Feature 15: Conversation Persistence}

Full conversation history is stored in SQLite (\kw{conversations} + \kw{messages} tables).
Streamlit sessions are ephemeral; \rapidsys{} adds server-side session caching
(via \kw{@st.cache\_resource}) keyed by a query-parameter session ID to survive page refreshes.

% ─────────────────────────────────────────────────────────────
\section{Feature 16: Text Preprocessing Pipeline}

Before chunking, raw extracted text passes through three stages:
\begin{enumerate}
  \item \textbf{OCR Artifact Cleaning} — fixes hyphenation artifacts, ligature substitutions (\texttt{fi} $\to$ \texttt{fi}), excessive spaces
  \item \textbf{Whitespace Normalization} — collapses 3+ newlines to 2, normalizes spaces
  \item \textbf{Heuristic Coreference Resolution} — replaces common pronouns with inferred referents using surrounding noun context (no ML model required)
\end{enumerate}

% ═══════════════════════════════════════════════════════════════
%  CHAPTER 5 — WORKFLOW ANALYSIS
% ═══════════════════════════════════════════════════════════════
\chapter{Workflow Analysis}

\section{Workflow 1: Document Upload}

\textbf{Starting Point:} User selects a file in the Streamlit UI.

\begin{enumerate}
  \item User clicks ``Upload File'' and selects a document
  \item Streamlit sends the file via multipart form POST to \kw{/upload}
  \item FastAPI receives the file; saves to \kw{uploads/} directory with UUID filename
  \item \kw{FileValidator.validate()} checks extension, MIME type, file size
  \item \kw{DocumentClassifier.classify()} analyzes content:
        \begin{itemize}
          \item Reads file extension
          \item Samples up to 2,000 words from text files
          \item Scores keyword frequency against domain vocabularies
          \item Returns \kw{ClassificationResult} (doc\_type, doc\_subtype, pipeline, confidence, stats)
        \end{itemize}
  \item \textbf{Branch A (tabular)} — \kw{TabularPipeline.ingest()} called:
        \begin{itemize}
          \item Reads file with pandas (auto-encoding detection)
          \item Cleans column names (sanitizes special characters)
          \item Loads into SQLite table in \kw{data/tabular\_uploads.db}
          \item Registers engine in \kw{CloudDatabaseService}; returns \kw{conn\_id}
        \end{itemize}
  \item \textbf{Branch B (narrative/RAG)} — \kw{RAGEngine.ingest\_document()} called:
        \begin{itemize}
          \item \kw{AutoConfigService} $\to$ optimal settings
          \item \kw{TextPreprocessor} $\to$ clean text
          \item \kw{LanguageDetector} $\to$ select embedding model
          \item \kw{ChunkingOptimizer} $\to$ best strategy $\to$ chunks
          \item \kw{EmbeddingManager} vectorizes each chunk
          \item Chunks stored in ChromaDB; BM25 index updated
          \item Knowledge graph extraction (async, optional)
        \end{itemize}
  \item \kw{DataCatalogService.register()} saves metadata + topics
  \item Document entry created in \kw{users.db}
  \item RBAC permissions set to \kw{private} (owner only, by default)
  \item Success response with \kw{doc\_id} returned to UI
\end{enumerate}

\textbf{Decision Points:}
\begin{itemize}
  \item File type not supported $\to$ HTTP 400 error
  \item Classification confidence $< 0.5$ $\to$ \kw{narrative/general/rag} fallback
  \item Tabular loading fails $\to$ error returned, no RAG fallback automatically
  \item ChromaDB unavailable $\to$ error
\end{itemize}

\textbf{End State:} Document indexed and immediately queryable.

% ─────────────────────────────────────────────────────────────
\section{Workflow 2: Query / Chat}

\textbf{Starting Point:} User types a question in the chat interface.

\begin{enumerate}
  \item User submits query via Streamlit chat input
  \item Streamlit POST to \kw{/query} with \kw{\{question, conversation\_id, provider, model\}}
  \item FastAPI validates JWT via \kw{SecurityService.get\_current\_user()}
  \item \kw{QueryDecomposer.decompose()} checks query complexity; splits if complex
  \item For each sub-query: \kw{orchestrator.\_invoke\_graph(sub\_query)} called
  \item LangGraph graph executes (classify $\to$ route $\to$ retrieve/generate $\to$ verify)
  \item If multiple sub-queries: \kw{QueryDecomposer.synthesize()} combines sub-answers
  \item RBAC filter applied to source citations
  \item Response includes: answer, sources, confidence level, conflict warnings
  \item Message saved to conversation in \kw{users.db}
  \item Response streamed back to Streamlit UI
\end{enumerate}

\textbf{Decision Points:}
\begin{itemize}
  \item \kw{general\_chat} detected $\to$ skip retrieval entirely
  \item No chunks retrieved $\to$ \kw{is\_unanswerable()} $\to$ trigger web search (if enabled)
  \item Low confidence after retry $\to$ partial answer with low-confidence flag
  \item DB query fails $\to$ error message returned
\end{itemize}

% ─────────────────────────────────────────────────────────────
\section{Workflow 3: Database Connection}

\textbf{Starting Point:} User wants to connect a live database.

\begin{enumerate}
  \item User enters host, port, database, username, password in ``Database'' settings
  \item POST to \kw{/database/connect} with credentials
  \item \kw{CloudDatabaseService.connect\_to\_postgres()} creates SQLAlchemy engine
  \item Test connection: lists tables to verify connectivity
  \item \kw{conn\_id} stored in memory; associated with user
  \item Schema context pre-fetched via \kw{get\_db\_schema\_context()}
  \item Database tables now queryable alongside uploaded documents
  \item On disconnect: \kw{close\_connection()} disposes SQLAlchemy engine
\end{enumerate}

% ─────────────────────────────────────────────────────────────
\section{Workflow 4: Cloud Storage Sync}

\textbf{Starting Point:} User selects files from Google Drive / S3 / etc.

\begin{enumerate}
  \item User authenticates via OAuth to cloud provider
  \item \rapidsys{} lists available files from cloud service
  \item User selects files to import
  \item Files downloaded to \kw{cloud\_cache/} directory
  \item Each file processed through normal upload pipeline (Workflow 1, steps 4--11)
  \item Cloud source URL stored in document metadata
\end{enumerate}

% ═══════════════════════════════════════════════════════════════
%  CHAPTER 6 — PROS AND CONS
% ═══════════════════════════════════════════════════════════════
\chapter{Pros and Cons}

\section{PROS — Strengths}

\begin{probox}[1. Zero-Configuration Intelligent Automation]
The most significant differentiator. Users upload a file and immediately get research-backed
optimal retrieval — no manual configuration. Derived from published academic benchmarks
(BEIR, QASPER, CUAD, PubMedQA, FinanceBench), giving the system a defensible and
reproducible quality baseline.
\end{probox}

\begin{probox}[2. Genuine Privacy-First Design]
\rapidsys{} is architectured to work 100\% offline: local LLMs via Ollama or LM Studio,
local embeddings via SentenceTransformers, local vector storage via ChromaDB, local SQLite
for all metadata. An organization can deploy it with zero internet connectivity.
\end{probox}

\begin{probox}[3. Sophisticated Dual-Pipeline Architecture]
Automatic routing of spreadsheet/tabular data to SQL and narrative text to RAG is
architecturally sound. Vector search on numeric data (prices, quantities, dates) is
fundamentally weak; SQL is the correct tool. \rapidsys{} makes this decision invisibly.
\end{probox}

\begin{probox}[4. Research-Backed Retrieval Parameters]
Unlike most RAG systems that use arbitrary chunk sizes (often 512 tokens everywhere),
\rapidsys{}'s \kw{AutoConfigService} cites specific benchmarks for each configuration.
For example, 950-word chunks for academic papers achieve +12\% F1 over 512-word chunks
on the QASPER benchmark.
\end{probox}

\begin{probox}[5. Multi-Tenancy with Fine-Grained RBAC]
Organizations, departments, groups, and document-level permissions are all implemented.
Document access checking is performed on every retrieval result. Audit logging provides
a compliance trail.
\end{probox}

\begin{probox}[6. Eleven Chunking Strategies]
Specialized strategies like \kw{table\_intact} (preserves markdown tables atomically),
\kw{code\_aware} (splits at function boundaries), and \kw{step\_aware} (splits at numbered
procedure steps) address specific retrieval failure modes not handled by standard RAG
frameworks.
\end{probox}

\begin{probox}[7. Hybrid Search with RRF Fusion]
Combining semantic and keyword search via Reciprocal Rank Fusion is a proven technique
that outperforms either alone. \rapidsys{} implements this properly using rank-based
aggregation rather than a simple weighted sum of incomparable score scales.
\end{probox}

\begin{probox}[8. Confidence Scoring with Retry Logic]
3-dimensional confidence scoring (context relevance + faithfulness + completeness) with
automatic retry is a meaningful quality gate. Most RAG systems return whatever the LLM
generates; \rapidsys{} has a mechanism to detect and remediate low-quality answers.
\end{probox}

\begin{probox}[9. Extensible Provider Architecture]
The abstract base class pattern for both \kw{LLMProvider} and \kw{EmbeddingProvider}
makes adding new providers straightforward. Singleton managers ensure consistent state
across the application.
\end{probox}

\begin{probox}[10. Comprehensive Cloud Storage Integration]
Native connectors for S3, Azure Blob, Google Drive, OneDrive, and Dropbox cover the
majority of enterprise storage use cases, reducing manual download-upload friction.
\end{probox}

\section{CONS — Weaknesses}

\begin{conbox}[1. Rate Limiting Not Production-Ready]
The rate limiter is in-memory using a Python dictionary. It resets on every server restart
and does not work across multiple server instances. In production (load-balanced), rate
limits are per-instance, not per-user globally. A proper production rate limiter requires
Redis or a distributed cache.
\end{conbox}

\begin{conbox}[2. BM25 Index Rebuild on Every Insert]
Every time a new document is indexed, the \kw{BM25Okapi} object is rebuilt from scratch
over the entire corpus. For a corpus with thousands of documents, this becomes increasingly
expensive and blocks the ingestion thread. The index should be made incremental or moved
to a proper inverted index.
\end{conbox}

\begin{conbox}[3. ChromaDB Collection Dimension Mismatch]
If a user switches embedding providers (or models) after documents are already indexed,
the new embeddings will have a different dimension than stored vectors. ChromaDB will
reject them or return garbage results. There is no migration path or detection of
this mismatch.
\end{conbox}

\begin{conbox}[4. No Token Budget Management]
When assembling chunks for the LLM prompt, the code uses a fixed \kw{top\_k} but does
not verify whether the total token count (chunks + question + system prompt + answer
budget) fits within the LLM's context window. Different LLMs have different limits
(4K, 8K, 128K). This can cause silent truncation or API errors.
\end{conbox}

\begin{conbox}[5. SQLite Concurrency Limitations]
SQLite is used for users, conversations, messages, documents, permissions, and the data
catalog — all in one file (\kw{users.db}). SQLite's write concurrency is limited (one
writer at a time). Under high user load, this creates a bottleneck.
\end{conbox}

\begin{conbox}[6. SQL Injection Risk from LLM-Generated SQL]
The \kw{execute\_query()} method uses SQLAlchemy's \kw{text()} with LLM-generated SQL.
If the LLM generates SQL with injection payloads (adversarial user input that leaks into
the prompt), it could execute arbitrary SQL. A proper implementation would restrict
to SELECT-only and validate against a schema whitelist.
\end{conbox}

\begin{conbox}[7. Conversation History Not Encrypted]
Messages stored in \kw{users.db} are plaintext. If the SQLite file is compromised, all
conversation content (which may include sensitive business information) is exposed. This
is a significant gap for a system designed to handle sensitive organizational data.
\end{conbox}

\begin{conbox}[8. Knowledge Graph Not Integrated with Confidence Loop]
The knowledge graph is built during ingestion but the \kw{graph\_query\_node} is a
separate code path from the verify/repair confidence loop. Graph queries cannot be
automatically retried with different strategies if they return low-confidence results.
\end{conbox}

\begin{conbox}[9. No Persistent Rate-Limit State]
Even if Redis were added, the current rate limit data structure does not persist state
across restarts. Persistent rate limiting requires a server-side store that outlives
the process.
\end{conbox}

\begin{conbox}[10. No UI for Configuration Management]
The \kw{AutoConfigService.update\_config()} method allows in-memory overrides, but there
is no UI for administrators to tune configurations. Config changes are not persisted
(restart resets them), making configuration management impractical without code changes.
\end{conbox}

% ═══════════════════════════════════════════════════════════════
%  CHAPTER 7 — REAL-WORLD SCENARIOS
% ═══════════════════════════════════════════════════════════════
\chapter{Real-World Problem Solving}

\section{Scenario 1: Law Firm Document Intelligence}

\subsection*{The Situation}
Hartley \& Associates is a mid-size law firm with 45 attorneys. Over 20 years they have
accumulated 80,000+ legal documents: contracts, case files, precedent analyses, and internal
policy documents stored across shared drives. New associates spend 2--4 hours per case doing
initial document review, searching for relevant precedents and clause examples. The firm
handles sensitive client information and cannot use public cloud AI tools.

\subsection*{Without RAPID}
\begin{itemize}
  \item Associates manually search through folder structures using Windows Search
  \item Relevant documents missed because search is keyword-only (misses synonyms, conceptual matches)
  \item Two attorneys independently research the same precedent, duplicating effort
  \item Junior associates lack the domain knowledge to know which search terms to use
  \item Average time for a contract review starting point: \textbf{3.5 hours}
\end{itemize}

\subsection*{With RAPID — Setup (done once by IT)}
\begin{enumerate}
  \item \rapidsys{} deployed on-premises (firm's own server)
  \item 80,000 documents bulk-uploaded via the cloud storage connector (Google Drive OAuth)
  \item \rapidsys{} auto-classifies each file:
        \begin{itemize}
          \item Legal contracts $\to$ \kw{legal/contract} (780-word chunks, hybrid search)
          \item Case analysis memos $\to$ \kw{narrative/policy}
          \item Billing spreadsheets $\to$ \kw{tabular/financial\_spreadsheet} (SQL pipeline)
        \end{itemize}
  \item All processing stays on-premises — client data never leaves the firm
\end{enumerate}

\subsection*{With RAPID — Daily Use}
An associate reviews a new software licensing agreement and needs to know how similar
indemnification clauses have been handled in past contracts.

\begin{enumerate}
  \item Opens \rapidsys{}, asks: ``Show me indemnification clauses from software licensing
        agreements where we represented the licensee, and how courts ruled on them in any
        associated case files.''
  \item \rapidsys{} decomposes: (a) indemnification clauses $\to$ legal contract RAG,
        (b) court rulings $\to$ case file RAG
  \item Retrieves relevant clause text from 8 historical contracts (hybrid search captures
        both ``indemnification'' keyword and semantic matches for ``hold harmless'')
  \item Retrieves relevant case outcome summaries
  \item ConfidenceScorer: CR\,=\,0.82, FA\,=\,0.79, CO\,=\,0.71 $\to$ overall\,=\,0.78 (high)
  \item Answer delivered with exact clause text, source document names, page references in \textbf{8 seconds}
\end{enumerate}

\subsection*{Outcome \& Impact}
\begin{itemize}
  \item Initial document review time: 3.5 hours $\to$ 20 minutes (\textbf{83\% reduction})
  \item Zero data leaves the firm's premises (compliance maintained)
  \item Junior associates produce research at near-senior-associate quality for precedent lookups
  \item After 6 months: firm estimates \textbf{1,200 associate-hours saved per month} across 45 attorneys
\end{itemize}

% ─────────────────────────────────────────────────────────────
\section{Scenario 2: Hospital Research Department Analytics}

\subsection*{The Situation}
Mercy General Hospital's research department conducts 40+ clinical trials simultaneously.
Each trial generates patient outcome spreadsheets, protocol documents, IRB approval letters,
published research papers, and regulatory correspondence. Department leadership needs to
answer funders' questions quickly — ``What was the 90-day readmission rate for Trial XR-44?''
or ``Summarize adverse events across all Drug Class B trials'' — but data is scattered across
Excel files and PDF reports.

\subsection*{Without RAPID}
\begin{itemize}
  \item Coordinator manually opens multiple Excel files, applies filters, copies numbers to new spreadsheets
  \item Summary reports require 4--6 hours of manual compilation
  \item Formatting errors occur when copying across spreadsheets
  \item Published literature must be read manually to cross-reference with internal findings
  \item If the lead coordinator is absent, institutional knowledge of where data lives is lost
\end{itemize}

\subsection*{With RAPID — Setup}
\begin{enumerate}
  \item Trial coordinators upload outcome spreadsheets (CSV/Excel) and protocol PDFs as they are created
  \item Spreadsheets auto-route to SQL pipeline — each becomes a queryable SQLite table
  \item PDFs auto-classify as \kw{medical/clinical\_document} $\to$ 500-word chunks, hybrid search
  \item Published papers auto-classify as \kw{academic/research\_paper} $\to$ 950-word section-aware chunks
  \item Data catalog maintains topic index with trial IDs, drug names, endpoints as keywords
\end{enumerate}

\subsection*{With RAPID — Funder Q\&A Session}

\textbf{Question 1:} ``What was the 90-day readmission rate for Trial XR-44?''
\begin{enumerate}
  \item DataCatalog finds the XR-44 outcome spreadsheet (topic match: ``XR-44'', ``readmission'')
  \item Routed to \kw{db\_agent} $\to$ generates SQL: \kw{SELECT AVG(readmission\_90d) FROM tbl\_xr44}
  \item Returns: ``The 90-day readmission rate for Trial XR-44 was 12.3\% (n\,=\,847 patients)''
  \item ConfidenceScorer: overall\,=\,0.91 (high confidence — exact SQL result)
\end{enumerate}

\textbf{Question 2:} ``Summarize adverse events across all Drug Class B trials''
\begin{enumerate}
  \item Query decomposed into separate sub-queries for each Drug Class B trial
  \item Parallel SQL queries across 6 trial tables for adverse event columns
  \item RAG query on protocol PDFs for adverse event definitions
  \item Synthesized: ``Across 6 Drug Class B trials (n\,=\,2,341 patients), most common adverse
        events were nausea (18.2\%), fatigue (14.7\%), and elevated liver enzymes (6.3\%)\ldots''
\end{enumerate}

\subsection*{Outcome \& Impact}
\begin{itemize}
  \item Funder Q\&A prep: 4--6 hours $\to$ \textbf{15 minutes} of conversational querying
  \item Zero errors from manual copy-paste between spreadsheets
  \item Research papers and trial data answered together in unified responses
  \item New coordinators can query institutional data from day one
  \item Estimated: \textbf{200 coordinator-hours saved per quarter} (equivalent to 1 FTE)
\end{itemize}

% ═══════════════════════════════════════════════════════════════
%  CHAPTER 8 — USE CASE DOCUMENTATION
% ═══════════════════════════════════════════════════════════════
\chapter{Use Case Documentation}

\section{UC-1: Upload and Query a Financial Spreadsheet}

\begin{table}[H]
\centering
\rowcolors{2}{rapidlight!40}{white}
\begin{tabularx}{\linewidth}{lX}
\toprule
\rowcolor{rapidblue!10}
\textbf{Field} & \textbf{Value} \\
\midrule
Actor & Finance analyst (role: \kw{user}) \\
Preconditions & Authenticated (valid JWT); org membership \\
\midrule
\multicolumn{2}{l}{\textbf{Main Flow}} \\
Step 1 & User selects \kw{quarterly\_sales.xlsx} in ``Upload'' tab \\
Step 2 & System validates file type and size \\
Step 3 & \kw{DocumentClassifier} detects tabular type (financial columns) \\
Step 4 & \kw{TabularPipeline} loads into SQLite; registers \kw{conn\_id} \\
Step 5 & \kw{DataCatalogService} registers with topics: [revenue, cost, margin, \ldots] \\
Step 6 & User asks: ``What was the highest revenue quarter in 2024?'' \\
Step 7 & Orchestrator routes to \kw{db\_agent\_node} \\
Step 8 & LLM generates SQL; query executed; result: ``Q3 2024, \$4.2M'' \\
Step 9 & Answer displayed with source citation \\
\midrule
Alternative Flows & File has encoding issues: tries UTF-8, latin-1, cp1252 automatically \\
Exception Flows & SQL generation failure: error returned; user prompted to rephrase \\
Postconditions & Spreadsheet queryable via NL indefinitely until user deletes it \\
Business Rules & Max 500,000 rows loaded; document private by default \\
\bottomrule
\end{tabularx}
\caption{UC-1: Upload and Query a Financial Spreadsheet}
\end{table}

\section{UC-2: Multi-Source Question}

\begin{table}[H]
\centering
\rowcolors{2}{rapidlight!40}{white}
\begin{tabularx}{\linewidth}{lX}
\toprule
\rowcolor{rapidblue!10}
\textbf{Field} & \textbf{Value} \\
\midrule
Actor & HR manager (role: \kw{manager}) \\
Preconditions & Policy PDF uploaded; employee database connected \\
\midrule
Step 1 & User asks: ``How many employees are on paternity leave and what does policy say about duration?'' \\
Step 2 & \kw{QueryDecomposer} splits into: (a) SQL query, (b) RAG query \\
Step 3 & DB agent queries employee table: \kw{SELECT COUNT(*) \ldots} \\
Step 4 & RAG retrieves policy PDF chunks about paternity leave \\
Step 5 & \kw{fusion\_node} merges count from DB + policy text from RAG \\
Step 6 & LLM synthesizes unified answer with dual-source citation \\
\midrule
Alternative Flow & No DB connected: only RAG answer, with note that headcount requires DB access \\
Postconditions & Answer with dual-source citation returned \\
\bottomrule
\end{tabularx}
\caption{UC-2: Multi-Source Question}
\end{table}

\section{UC-3: Admin Creates a New User}

\textbf{Actor:} Administrator (role: \kw{admin})

\textbf{Main Flow:}
\begin{enumerate}
  \item Admin calls \kw{POST /admin/users} with \kw{\{username, password, role, org\_id, department, groups\}}
  \item Password validated (min 8 chars, complexity requirements)
  \item User created in \kw{users.db}; groups pre-populated if specified
  \item Admin shares documents with the new user via \kw{PUT /documents/\{doc\_id\}/permissions}
\end{enumerate}

\textbf{Exception Flows:} Username already exists $\to$ HTTP 400; Weak password $\to$ HTTP 400; Invalid role $\to$ HTTP 400.

\section{UC-4: Connect to External Database}

\textbf{Actor:} Data analyst (role: \kw{manager})

\textbf{Main Flow:}
\begin{enumerate}
  \item User opens ``Databases'' settings panel; enters host, port, database, username, password
  \item POST to \kw{/database/connect}
  \item \kw{CloudDatabaseService.connect\_to\_postgres()} creates SQLAlchemy engine
  \item Test connection; schema pre-fetched via \kw{get\_db\_schema\_context()}
  \item Connection available immediately in chat interface
\end{enumerate}

\section{UC-5: Query Returns Low Confidence — Retry Loop}

\textbf{Main Flow:}
\begin{enumerate}
  \item User asks a highly specific technical question not well-covered in documents
  \item Orchestrator retrieves chunks; ConfidenceScorer: CR\,=\,0.25, FA\,=\,0.41, CO\,=\,0.30 $\to$ overall\,=\,0.34 (below 0.40)
  \item \kw{repair\_node} activated: HyDE generates a hypothetical passage; used as new query vector
  \item Re-retrieval with HyDE vector finds more relevant chunks
  \item Re-evaluation: overall\,=\,0.58 $\to$ ``medium'' verdict
  \item Answer returned with ``\(\triangle\) Low-confidence answer'' flag and suggestion to check source directly
\end{enumerate}

% ═══════════════════════════════════════════════════════════════
%  CHAPTER 9 — ORCHESTRATION & ARCHITECTURE
% ═══════════════════════════════════════════════════════════════
\chapter{Orchestration \& Architecture}

\section{LangGraph State Machine}

The core orchestration is a \textbf{directed state graph} implemented with LangGraph.
The graph state is a TypedDict:

\begin{codebox}[AgentState TypedDict]
class AgentState(TypedDict):
    query: str
    intent: str
    messages: List[BaseMessage]
    rag_context: List[str]
    db_results: Optional[str]
    final_answer: str
    confidence: float
    retry_count: int
    sources: List[str]
\end{codebox}

\subsection*{Conditional Edges}
\begin{description}
  \item[\kw{classify\_node}:] routes based on intent label to the appropriate agent node
  \item[\kw{verify\_node}:] score $\geq 0.65 \to$ END; score $< 0.40$ and retry $< 1 \to$ \kw{repair\_node}; else $\to$ \kw{partial\_deliver\_node}
\end{description}

\section{Service Singleton Pattern}

Both \kw{LLMManager} and \kw{EmbeddingManager} implement the \textbf{Singleton pattern}
using \kw{\_\_new\_\_} override. This ensures:
\begin{itemize}
  \item A single model instance loaded in memory (avoids re-loading multi-GB transformer models)
  \item Consistent active provider/model state across all requests
  \item Thread safety for provider selection
\end{itemize}

\section{Pipeline Orchestration}

\begin{codebox}[Upload Orchestration Flow]
File Upload
    |
    |-- FileValidator (gate)
    |
    |-- DocumentClassifier (classify)
    |       |
    |       |-- tabular --> TabularPipeline
    |       |       +-- SQLite --> CloudDatabaseService.register()
    |       |
    |       +-- text/code --> RAGEngine
    |               |-- AutoConfigService  (settings)
    |               |-- TextPreprocessor  (clean)
    |               |-- LanguageDetector  (embedding model)
    |               |-- ChunkingOptimizer (strategy + chunk)
    |               |-- EmbeddingManager  (vectorize)
    |               |-- ChromaDB          (store vectors)
    |               |-- FullTextSearchEngine (BM25)
    |               +-- KnowledgeGraphBuilder (async)
    |
    +-- DataCatalogService.register() (metadata)
\end{codebox}

\section{State Management}

Application state is managed at three levels:
\begin{enumerate}
  \item \textbf{Request level}: FastAPI dependency injection (current\_user, DB connections)
  \item \textbf{Session level}: Streamlit \kw{session\_state} + server-side cache keyed by session ID
  \item \textbf{Persistent level}: SQLite (users, documents, conversations), ChromaDB (vectors),
        JSON files (BM25 index, knowledge graphs), SQLite (tabular data)
\end{enumerate}

\section{Data Storage Layout}

\begin{codebox}[File System Layout]
data/
|-- users.db              # SQLite: users, orgs, groups, docs,
|                         #         conversations, data_catalog
|-- tabular_uploads.db    # SQLite: uploaded spreadsheet data
|-- chroma/               # ChromaDB vector store (persistent)
|-- fulltext_index.db     # SQLite: BM25 search metadata
|-- search/bm25_index.json  # BM25 corpus
|-- knowledge_graph/      # NetworkX graphs (per-user JSON)
|-- traces.db             # Query traces for debugging
|-- audit.log             # Rotating security audit log
|-- .jwt_secret           # Auto-generated JWT key (chmod 600)
+-- .encryption_key       # Encryption service key
\end{codebox}

% ═══════════════════════════════════════════════════════════════
%  CHAPTER 10 — TECHNOLOGY STACK
% ═══════════════════════════════════════════════════════════════
\chapter{Technology Stack \& Libraries}

\section{Backend Framework}

\begin{table}[H]
\centering
\rowcolors{2}{rapidlight!40}{white}
\begin{tabularx}{\linewidth}{llX}
\toprule
\rowcolor{rapidblue!10}
\textbf{Technology} & \textbf{Version} & \textbf{Purpose} \\
\midrule
FastAPI & 0.104.0 & REST API framework; async support; Pydantic validation \\
Uvicorn & 0.24.0 & ASGI server for FastAPI \\
Gunicorn & 21.0+ & Production process manager for Uvicorn workers \\
\bottomrule
\end{tabularx}
\caption{Backend framework}
\end{table}

\section{Frontend}

\begin{table}[H]
\centering
\rowcolors{2}{rapidlight!40}{white}
\begin{tabularx}{\linewidth}{llX}
\toprule
\rowcolor{rapidblue!10}
\textbf{Technology} & \textbf{Version} & \textbf{Purpose} \\
\midrule
Streamlit & 1.30.0 & Web UI; chat interface, file upload, settings panels \\
\bottomrule
\end{tabularx}
\caption{Frontend framework}
\end{table}

\section{AI / ML Layer}

\begin{table}[H]
\centering
\rowcolors{2}{rapidlight!40}{white}
\begin{tabularx}{\linewidth}{llX}
\toprule
\rowcolor{rapidblue!10}
\textbf{Library} & \textbf{Version} & \textbf{Purpose} \\
\midrule
LangChain OpenAI & 0.1.0 & LangChain adapter for OpenAI models \\
LangChain Anthropic & 0.1.0 & LangChain adapter for Claude models \\
LangChain Community & 0.0.20+ & Ollama LLM adapter \\
LangGraph & 0.0.40+ & State machine for multi-agent orchestration \\
OpenAI SDK & 1.109.1 & Direct OpenAI API calls, embeddings \\
Anthropic SDK & 0.16.0+ & Direct Anthropic API calls \\
sentence-transformers & 2.2.0+ & Local embedding models (all-MiniLM-L6-v2, multilingual-e5) \\
\bottomrule
\end{tabularx}
\caption{AI / ML libraries}
\end{table}

\section{Vector Store \& Search}

\begin{table}[H]
\centering
\rowcolors{2}{rapidlight!40}{white}
\begin{tabularx}{\linewidth}{llX}
\toprule
\rowcolor{rapidblue!10}
\textbf{Library} & \textbf{Version} & \textbf{Purpose} \\
\midrule
ChromaDB & 0.4.22 & Vector database for semantic search \\
rank-bm25 & 0.2.2+ & BM25 Okapi keyword search \\
\bottomrule
\end{tabularx}
\caption{Vector store and keyword search}
\end{table}

\section{Document Processing}

\begin{table}[H]
\centering
\rowcolors{2}{rapidlight!40}{white}
\begin{tabularx}{\linewidth}{llX}
\toprule
\rowcolor{rapidblue!10}
\textbf{Library} & \textbf{Version} & \textbf{Purpose} \\
\midrule
pypdf & 4.0.0 & PDF text extraction \\
pdfplumber & 0.10.0 & Advanced PDF extraction (tables, layout) \\
python-docx & 1.1.0 & Word document parsing \\
openpyxl & 3.1.0 & Excel file reading \\
beautifulsoup4 & 4.12.0 & HTML parsing \\
chardet & 5.2.0+ & Character encoding detection \\
langdetect & 1.0.9+ & Language detection \\
pyarrow & 14.0.0+ & Parquet file reading \\
\bottomrule
\end{tabularx}
\caption{Document processing libraries}
\end{table}

\section{Database Layer}

\begin{table}[H]
\centering
\rowcolors{2}{rapidlight!40}{white}
\begin{tabularx}{\linewidth}{llX}
\toprule
\rowcolor{rapidblue!10}
\textbf{Technology} & \textbf{Version} & \textbf{Purpose} \\
\midrule
SQLAlchemy & 2.0.36+ & ORM for PostgreSQL, MySQL, Snowflake, BigQuery \\
aiosqlite & 0.19.0 & Async SQLite access \\
sqlite3 & stdlib & User DB, catalog, conversations, tabular data \\
pymysql & 1.1.0 & MySQL driver \\
snowflake-connector-python & 3.7.1 & Snowflake connectivity \\
google-cloud-bigquery & 3.13.0 & BigQuery connectivity \\
\bottomrule
\end{tabularx}
\caption{Database layer}
\end{table}

\section{Data Processing}

\begin{table}[H]
\centering
\rowcolors{2}{rapidlight!40}{white}
\begin{tabularx}{\linewidth}{llX}
\toprule
\rowcolor{rapidblue!10}
\textbf{Library} & \textbf{Version} & \textbf{Purpose} \\
\midrule
pandas & latest & DataFrame operations, CSV/Excel/Parquet loading \\
numpy & $< 2$ & Numerical operations (version capped for compatibility) \\
networkx & 3.0+ & Knowledge graph (directed graph, graph algorithms) \\
\bottomrule
\end{tabularx}
\caption{Data processing libraries}
\end{table}

\section{Security}

\begin{table}[H]
\centering
\rowcolors{2}{rapidlight!40}{white}
\begin{tabularx}{\linewidth}{llX}
\toprule
\rowcolor{rapidblue!10}
\textbf{Library} & \textbf{Version} & \textbf{Purpose} \\
\midrule
PyJWT & 2.8.0 & JWT token creation and verification \\
bcrypt & 4.1.2 & Password hashing \\
cryptography & 41.0.0+ & Encryption service for file contents at rest \\
msal & 1.24.0+ & Microsoft OAuth / Azure AD authentication \\
\bottomrule
\end{tabularx}
\caption{Security libraries}
\end{table}

\section{Cloud Storage \& Integration}

\begin{table}[H]
\centering
\rowcolors{2}{rapidlight!40}{white}
\begin{tabularx}{\linewidth}{llX}
\toprule
\rowcolor{rapidblue!10}
\textbf{Library} & \textbf{Version} & \textbf{Purpose} \\
\midrule
boto3 & 1.34.0+ & AWS S3 operations \\
azure-storage-blob & 12.19.0+ & Azure Blob Storage \\
google-api-python-client & 2.100.0+ & Google Drive API \\
google-auth-oauthlib & 1.1.0+ & Google OAuth flow \\
dropbox & 11.36.0+ & Dropbox API \\
\bottomrule
\end{tabularx}
\caption{Cloud storage connector libraries}
\end{table}

\section{Infrastructure \& Utilities}

\begin{table}[H]
\centering
\rowcolors{2}{rapidlight!40}{white}
\begin{tabularx}{\linewidth}{llX}
\toprule
\rowcolor{rapidblue!10}
\textbf{Library} & \textbf{Version} & \textbf{Purpose} \\
\midrule
redis & 5.0.0 & Query result caching (optional) \\
httpx & 0.25.0 & Async HTTP for web search \\
requests & 2.31.0 & Sync HTTP for provider availability checks \\
Jinja2 & 3.1.0 & Template rendering \\
python-multipart & 0.0.6 & Multipart form data (file uploads) \\
pytest & 7.4.0 & Test framework \\
reportlab & 4.0.0+ & PDF report generation \\
\bottomrule
\end{tabularx}
\caption{Infrastructure and utility libraries}
\end{table}

% ═══════════════════════════════════════════════════════════════
%  CHAPTER 11 — APPENDIX: GLOSSARY
% ═══════════════════════════════════════════════════════════════
\chapter{Appendix — Glossary}

\begin{longtable}{>{\bfseries}p{3.2cm}p{11.3cm}}
\toprule
\rowcolor{rapidblue!10}
\textbf{Term} & \textbf{Definition} \\
\midrule
\endhead
RAG & Retrieval-Augmented Generation. AI technique where relevant documents are retrieved and included in the prompt to ground LLM answers in factual context. \\
BM25 & Okapi BM25. A probabilistic ranking function for keyword-based document retrieval. Outperforms TF-IDF for most search tasks. \\
ChromaDB & Open-source vector database. Stores document embeddings and enables nearest-neighbor similarity search. \\
Embedding & A numerical vector representation of text. Semantically similar texts produce similar vectors (close together in vector space). \\
Chunk & A portion of a document. Large documents are split into chunks so each fits in the LLM context and can be individually retrieved. \\
RRF & Reciprocal Rank Fusion. A rank aggregation method that combines results from multiple retrieval systems without requiring score normalization. \\
LangGraph & A library for building stateful, multi-step AI workflows as directed graphs. Each node is a processing step; edges define control flow. \\
CRAG & Corrective RAG. An architecture pattern where retrieved documents are evaluated and corrected/augmented before generation if quality is insufficient. \\
HyDE & Hypothetical Document Embedding. The LLM generates a hypothetical answer, which is embedded and used as the retrieval query. Improves recall for vague queries. \\
RBAC & Role-Based Access Control. Authorization scheme where permissions are assigned to roles rather than individual users. \\
JWT & JSON Web Token. A compact, self-contained token for transmitting authentication information between parties. \\
NL-to-SQL & Natural Language to SQL. Converting plain English questions to SQL queries automatically using an LLM. \\
SentenceTransformers & Python library providing pre-trained transformer models optimized for generating sentence and document embeddings. \\
Ollama & Open-source tool for running large language models locally on consumer hardware. \\
Speculative RAG & A two-LLM pattern: a fast small model drafts the answer, a larger model verifies it. Accepted drafts skip verification. \\
QASPER & A dataset benchmark for question answering over scientific papers. Used to evaluate academic document retrieval strategies. \\
CUAD & Contract Understanding Atticus Dataset. A benchmark for evaluating AI on legal contract analysis. \\
FinanceBench & A benchmark dataset for evaluating financial document question answering performance. \\
BEIR & Benchmarking IR. A heterogeneous benchmark for information retrieval systems across diverse domains. \\
NER & Named Entity Recognition. Identifying and classifying named entities (persons, organizations, locations) in text. \\
NetworkX & Python library for creating, analyzing, and visualizing complex graphs and networks. \\
OCR & Optical Character Recognition. Technology that converts scanned images of text into machine-readable characters. \\
ASGI & Asynchronous Server Gateway Interface. Python specification for async web servers and applications. \\
ORM & Object-Relational Mapping. Technique that lets you query and manipulate a database using an object-oriented paradigm. \\
\bottomrule
\caption{Glossary of Technical Terms}
\end{longtable}

% ═══════════════════════════════════════════════════════════════
%  BACK MATTER
% ═══════════════════════════════════════════════════════════════
\cleardoublepage
\begin{center}
  \vspace*{8cm}
  {\Large\sffamily\bfseries\color{rapidblue} End of Report}\\[1cm]
  {\normalsize\sffamily\color{rapidgray}
    RAPID — RAG Application for Private Instant Deployment\\
    Comprehensive Technical Analysis\\[0.5cm]
    Analysis Date: February 27, 2026\\
    Document Version: 1.0
  }
\end{center}

\end{document}
